% Source: https://artofproblemsolving.com/community/c3h1929832_math_problems
% Stable identity: aops-topic-1929832-q50
The semiperimeter is $s=21$ and the area is $\sqrt{21\cdot7\cdot6\cdot8}=84$, so the inradius is $r=\dfrac{84}{21}=4$. The tangent lengths from the vertices to the incircle are $s-BC=7$ from $A$, $s-CA=6$ from $B$ and $s-AB=8$ from $C$. The radius to a point of tangency is perpendicular to the side, so by the Pythagorean theorem
\[
AI^2=7^2+4^2=65,\qquad BI^2=6^2+4^2=52,\qquad CI^2=8^2+4^2=80 .
\]

\begin{asy}
usepackage("amsmath");
size(6cm);
defaultpen(fontsize(10pt));
pair B=(0,0), C=(14,0), A=(5,12), I=(6,4);
pair Ta=(6,0), Tb=C+8*unit(A-C), Tc=6*unit(A);
draw(circle(I,4),gray);
draw(A--B--C--cycle,linewidth(1));
draw(A--I,blue+linewidth(0.9)); draw(B--I,blue+linewidth(0.9)); draw(C--I,blue+linewidth(0.9));
draw(I--Ta,dashed);
dot(I);dot(Ta);dot(Tb);dot(Tc);
label("$A$",A,N); label("$B$",B,SW); label("$C$",C,SE); label("$I$",I,NE);
label("$6$",(B+Ta)/2,S); label("$8$",(Ta+C)/2,S);
label("$6$",(B+Tc)/2,W); label("$7$",(Tc+A)/2,W);
label("$8$",(C+Tb)/2,NE); label("$7$",(Tb+A)/2,NE);
label("$4$",(I+Ta)/2,E);
label("\mbox{}",(min(currentpicture,user=true)+max(currentpicture,user=true))/2,black); // keeps the text after the figure black when LaTeX typesets the labels inline
\end{asy}

For a triangle with sides $x,y,z$ and area $K$, Heron's formula can be written as
\[
16K^2=2\left(x^2y^2+y^2z^2+z^2x^2\right)-\left(x^4+y^4+z^4\right).
\]
With $x^2=65$, $y^2=52$, $z^2=80$ (these satisfy the triangle inequality, e.g.\ $\sqrt{52}+\sqrt{65}>\sqrt{80}$):
\[
16K^2=2(3380+4160+5200)-(4225+2704+6400)=25480-13329=12151,
\]
so $K^2=\boxed{\dfrac{12151}{16}}$.
